\documentclass[12pt,a4paper]{article}
\usepackage[utf8]{inputenc}
\usepackage[T1]{fontenc}
\usepackage{geometry}
\usepackage{amsmath}
\geometry{margin=2.2cm}
\setlength{\parindent}{0pt}
\setlength{\parskip}{0.55em}
\title{Synthetic Validation Guide\\Using \texttt{run\_synthetic\_validation.py} in the current project}
\author{Krogh GUI project note}
\date{23 April 2026}

\begin{document}
\sloppy
\maketitle
\tableofcontents

\section{Purpose}
The synthetic validation workflow checks forward model and inverse reconstruction behavior on fictive but physiologically plausible cases. It is intended to detect numerical regressions, weak identifiability regions, and trend violations before relying on real-world datasets.

\section{How to run}
From project root:
\begin{verbatim}
source .venv/bin/activate
python run_synthetic_validation.py
\end{verbatim}

The script creates \texttt{Diagnostic reports/} if needed and prints progress lines.

\section{What the script executes internally}
\texttt{run\_synthetic\_validation.py} calls:
\begin{center}
\texttt{krogh\_app.validation.run\_and\_save\_default\_validation(output\_dir)}
\end{center}

Current default sequence:
\begin{enumerate}
    \item baseline synthetic suite,
    \item monotonic trend checks,
    \item compact noise-robustness checks,
    \item CSV and JSON export for all sections.
\end{enumerate}

\section{Output files and locations}
All files are written to:
\begin{verbatim}
Diagnostic reports/
\end{verbatim}

Generated by default run:
\begin{itemize}
    \item \texttt{synthetic\_validation\_report.csv}
    \item \texttt{synthetic\_validation\_report.json}
    \item \texttt{synthetic\_validation\_trend\_checks.csv}
    \item \texttt{synthetic\_validation\_trend\_checks.json}
    \item \texttt{synthetic\_validation\_robustness.csv}
    \item \texttt{synthetic\_validation\_robustness.json}
\end{itemize}

\section{How to read the baseline report}
The baseline report summarizes one row per synthetic case (plus summary section in JSON). Typical fields include status, error terms, fit diagnostics, and reconstruction quality indicators.

Key interpretation principles:
\begin{itemize}
    \item \texttt{pass} means current acceptance criteria were met.
    \item \texttt{review} means inspect the case; this often indicates parameter non-uniqueness rather than software failure.
    \item sensor and venous errors should be interpreted together with uncertainty/identifiability signals.
\end{itemize}

\section{How to read trend checks}
Trend checks verify directional plausibility. In current implementation, examples include checking whether sensor-like tissue PO$_2$ increases when:
\begin{itemize}
    \item perfusion is increased,
    \item inlet PO$_2$ is increased.
\end{itemize}

Monotonicity flags in summary provide a fast physiology-consistency sanity check.

\section{How to read robustness checks}
Robustness checks perturb input targets by small noise levels and evaluate how much fit quality degrades.

Interpretation:
\begin{itemize}
    \item small degradation is expected,
    \item large, unstable jumps suggest fragile operating regions or weak identifiability.
\end{itemize}

\section{Progress behavior and runtime expectations}
Normal behavior:
\begin{itemize}
    \item progress lines for baseline, trend, robustness stages,
    \item visible case-by-case updates,
    \item output files appear and update in \texttt{Diagnostic reports/}.
\end{itemize}

Intervention is usually only needed if output stops for an unusually long period and there is no active CPU progress.

\section{Mathematical context of what is being validated}
Synthetic validation probes the same model components documented in the main manual:
\begin{itemize}
    \item effective $P_{50}$ shift,
    \item Hill saturation and differential capacity,
    \item Krogh transport plus tissue diffusion coupling,
    \item inverse fitting of perfusion, mitoP50, and metabolic factor,
    \item uncertainty and identifiability summaries.
\end{itemize}

Therefore, validation outcomes are directly relevant for the trustworthiness of both forward simulations and diagnostic reconstruction interpretations.

\section{Recommended practical workflow}
\begin{enumerate}
    \item Run synthetic validation after solver/fitting changes.
    \item Review baseline pass/review cases.
    \item Confirm trend checks remain monotonic.
    \item Review robustness degradation levels.
    \item Only then proceed to broader benchmark runs or manuscript figures.
\end{enumerate}

\section{Relation to reconstruction benchmark}
Synthetic validation and reconstruction benchmark are complementary:
\begin{itemize}
    \item synthetic validation focuses on model-behavior quality and robustness,
    \item reconstruction benchmark focuses on runtime and fit quality across representative benchmark cases (including optimized vs legacy-grid comparison).
\end{itemize}

Benchmark command:
\begin{verbatim}
python run_reconstruction_benchmark.py
\end{verbatim}

\section{Conclusion}
The synthetic validation tool is a reproducible quality-control layer for this project. It does not replace real-world validation, but it is essential for catching regressions, documenting stability, and understanding where inverse fits are strong versus only weakly identifiable.

\end{document}
